\section{Introduction}
\label{sec:introduction}

Intracranial hemorrhage (ICH) is among the most severe neurological
emergencies, carrying a 30-day mortality rate of approximately
40--50\,\% and leaving a large proportion of survivors with permanent
disability~\cite{feigin2013global,lovelock2010incidence}.
Rapid diagnosis is decisive: each hour of delayed treatment is
associated with meaningfully worse functional outcomes, and the
recommended window for administering reversal agents in anticoagulation-related
ICH is under four hours from onset~\cite{hemphill2015ahaguideline}.
Non-contrast head computed tomography (NCCT) is the clinical gold standard
for initial ICH evaluation because of its availability, speed, and high
sensitivity to acute blood~\cite{broderick2007ich}.
Despite this, the interpretation of head CT scans requires specialist
radiologist expertise that is unevenly distributed across
institutions and time shifts, motivating automated decision-support systems
that can flag suspicious scans and classify hemorrhage subtype.

The 2019 RSNA Intracranial Hemorrhage Detection challenge~\cite{rsna2019}
provided a large public benchmark---over 25\,000 CT studies labelled by
a panel of radiologists---and catalysed rapid progress in deep learning for
ICH detection.
The first-place solution~\cite{wang2021deeplearning} demonstrated that a
two-stage pipeline combining a per-slice convolutional neural network (CNN)
with a study-level bidirectional recurrent sequence model outperforms
single-image classifiers, because ICH often spans multiple adjacent slices
and the inter-slice context carries diagnostic information not available from
any single image alone.
A subsequent systematic review~\cite{cho2021systematic} confirmed that
ensemble methods and attention mechanisms consistently improve sensitivity
and specificity across diverse ICH datasets.

Building on these observations, this paper reimplements and extends the
first-place RSNA 2019 pipeline on the publicly available CQ500
dataset~\cite{chilamkurthy2018cq500}, which contains 491 CT studies
with per-study consensus labels from three radiologists.
Our contributions are:

\begin{itemize}
    \item A three-backbone ensemble (DenseNet-121, DenseNet-169,
          SE-ResNeXt-101) for Stage-1 per-slice ICH classification with
          six output classes, trained on three-channel adjacent-slice images
          generated from CT windowing.
    \item A Stage-2 bidirectional GRU sequence model that refines per-slice
          predictions using intra-study context, augmented with the \emph{slice
          thickness} DICOM metadata as an additional input channel---a
          feature described in the original paper but absent from most
          reproductions.
    \item Bootstrap-based 95\,\% confidence interval estimation for all
          reported AUC, sensitivity, and specificity values, enabling
          statistically meaningful comparison with published benchmarks.
    \item An open-source, reproducible codebase with a
          complete DVC-managed data pipeline, MLflow experiment tracking,
          and a REST inference API served in Docker.
\end{itemize}

\begin{figure}[t]
    \centering
    \includegraphics[width=\linewidth]{../../reports/figures/gradcam.png}
    \caption{Grad-CAM activation maps overlaid on CT slices for the
             \emph{any hemorrhage} class after Stage-1 inference.
             Warmer colours indicate regions most influential for the
             positive prediction.  The model consistently attends to
             hyperdense (bright) regions characteristic of acute blood.}
    \label{fig:gradcam_overview}
\end{figure}

The remainder of this paper is organised as follows.
Section~\ref{sec:related_work} reviews related work on deep learning for
ICH and CT image analysis.
Section~\ref{sec:methodology} describes the preprocessing pipeline,
network architectures, and training procedure.
Section~\ref{sec:experiments} details the experimental setup, dataset,
and evaluation metrics.
Section~\ref{sec:results} presents quantitative results and qualitative
analysis.
Section~\ref{sec:conclusion} concludes and outlines directions for
future work.
